%! TEX root = main.tex  

\chapter{Response of 1-DOF Systems}

\section{Free Vibration}

\subsection{Undamped Case}

Consider the free vibration of an undamped single degree of freedom system:
\[
m \ddot{x} + kx = 0
\]
Assume a solution of the form:
\[
x(t) = X e^{\lambda t} \quad \Rightarrow \quad \dot{x}(t) = \lambda X e^{\lambda t}, \quad \ddot{x}(t) = \lambda^2 X e^{\lambda t}
\]
Substituting into the equation of motion:
\[
m \lambda^2 X e^{\lambda t} + k X e^{\lambda t} = 0 \quad \Rightarrow \quad (m \lambda^2 + k) X e^{\lambda t} = 0
\]
Since \( e^{\lambda t} \neq 0 \) and \( X \neq 0 \) (non-trivial solution), the characteristic equation is:
\[
m \lambda^2 + k = 0 \quad \Rightarrow \quad \lambda_{1,2} = \pm j \sqrt{\frac{k}{m}} = \pm j \omega_0
\]
Thus, the general solution is:
\[
x(t) = X_1 e^{j \omega_0 t} + X_2 e^{-j \omega_0 t}
\]
Since \( x(t) \in \mathbb{R} \), let:
\[
X_1 = \frac{A - jB}{2}, \quad X_2 = \frac{A + jB}{2}
\]
So that:
\begin{align*}
x(t) &= X_1 \left[ \cos(\omega_0 t) + j \sin(\omega_0 t) \right] + X_2 \left[ \cos(\omega_0 t) - j \sin(\omega_0 t) \right] \\
&= (X_1 + X_2) \cos(\omega_0 t) + j (X_1 - X_2) \sin(\omega_0 t) \\
&= A \cos(\omega_0 t) + B \sin(\omega_0 t)
\end{align*}

Alternatively, define:
\[
A = C \cos \phi, \quad B = C \sin \phi \quad \Rightarrow \quad C = \sqrt{A^2 + B^2}, \quad \phi = \arctan\left( \frac{B}{A} \right)
\]
so that:
\[
x(t) = C \cos(\omega_0 t - \phi)
\]
This is the standard harmonic oscillator solution.

\subsection{Damped Case}

Now consider the damped system:
\[
m \ddot{x} + c \dot{x} + k x = 0
\]
Assume again \( x(t) = X e^{\lambda t} \), \( \dot{x}(t) = \lambda X e^{\lambda t} \), \( \ddot{x}(t) = \lambda^2 X e^{\lambda t} \), then substituting:
\[
m \lambda^2 X e^{\lambda t} + c \lambda X e^{\lambda t} + k X e^{\lambda t} = 0
\quad \Rightarrow \quad 
(m \lambda^2 + c \lambda + k) X e^{\lambda t} = 0
\]
As before, since \( e^{\lambda t} \neq 0 \) and \( X \neq 0 \) (non-trivial solution), the characteristic equation is:
\[
m \lambda^2 + c \lambda + k = 0 
\quad \Rightarrow \quad
\lambda_{1,2} = \frac{-c \pm \sqrt{c^2 - 4mk}}{2m} = -\alpha \pm \sqrt{\alpha^2 - \omega_0^2}
\]
where \( \alpha = \frac{c}{2m} \) and \( \omega_0 = \sqrt{\frac{k}{m}} \).

\paragraph{Critically Damped Case} If \( \Delta = 0 \), i.e., \( \alpha = \omega_0 \), the critical damping value can be computed by requiring the discriminant of the characteristic equation to be zero:
\[
\Delta = c^2 - 4mk = 0 \quad \Rightarrow \quad c_{\text{cr}} = 2 \sqrt{km}
\]
This yields:
\[
\alpha = \frac{c_{\text{cr}}}{2m} = \frac{1}{2m} \cdot 2\sqrt{km} = \sqrt{\frac{k}{m}} = \omega_0
\]
As expected, in the critically damped case \( \alpha = \omega_0 \), leading to a double coincident real root:
\[
\lambda_{1,2} = -\alpha
\]
The solution becomes:
\[
x(t) = X_1 e^{-\alpha t} + X_2 t e^{-\alpha t}
\]
There is no oscillation. 

\paragraph{Underdamped Case} If \( \Delta < 0 \), i.e., \( \alpha < \omega_0 \), or equivalently \( c < c_{\text{cr}} \), then it yields two complex conjugate solutions:
\[
\lambda_{1,2} = -\alpha \pm j \omega_d \quad \text{where} \quad \omega_d = \sqrt{\omega_0^2 - \alpha^2}
\]
The solution becomes, using the same decomposition as in the undamped case (complex exponential, trigonometric, and phase form):
\[
x(t) = e^{-\alpha t} (X_1 e^{j \omega_d t} + X_2 e^{-j \omega_d t}) 
= e^{-\alpha t} (A \cos(\omega_d t) + B \sin(\omega_d t)) 
= C e^{-\alpha t} \cos(\omega_d t - \phi)
\]
This is the damped harmonic oscillator response.

\paragraph{Overdamped Case} If \( \Delta > 0 \), i.e., \( \alpha > \omega_0 \), or \( c > c_{\text{cr}} \), then two negative real solutions can be computed:
\[
\lambda_{1,2} = -\alpha \pm \beta, \quad \text{with } \beta = \sqrt{\alpha^2 - \omega_0^2}
\]
Since \( \beta < \alpha \), both roots are real and negative:
\[
\lambda_{1,2} \in \mathbb{R}^- \quad\Rightarrow \quad \lambda_{1,2} = -\alpha_{1,2}
\]
The solution becomes:
\[
x(t) = X_1 e^{\lambda_1 t} + X_2 e^{\lambda_2 t} = X_1 e^{-\alpha_1 t} + X_2 e^{-\alpha_2 t}
\]
No oscillation occurs in this case.

\paragraph{Definition of the Damping Ratio} Define the dimensionless damping ratio:
\[
h = \frac{\alpha}{\omega_0} = \frac{c}{2 \sqrt{km}}
\quad \Rightarrow \quad 
\lambda_{1,2} = -\alpha \pm j \omega_0 \sqrt{1 - h^2}
\]
This compact form helps to distinguish between underdamped (\( h < 1 \)), critically damped (\( h = 1 \)), and overdamped (\( h > 1 \)) systems.

\subsection{Determination of Amplitude Parameters from Initial Conditions}

Given the underdamped solution in trigonometric form:
\[
x(t) = e^{-\alpha t} (A \cos(\omega_d t) + B \sin(\omega_d t))
\]
and initial conditions \( x(0) = x_0 \), \( \dot{x}(0) = v_0 \), determine the constants \( A \) and \( B \). At \( t = 0 \), the solution becomes:
\[
x(0) = A = x_0
\]
Taking the derivative:
\[
\dot{x}(t) = -\alpha e^{-\alpha t} (A \cos(\omega_d t) + B \sin(\omega_d t)) + e^{-\alpha t} (-A \omega_d \sin(\omega_d t) + B \omega_d \cos(\omega_d t))
\]
Evaluating at \( t = 0 \):
\[
\dot{x}(0) = -\alpha A + B \omega_d = v_0
\quad \Rightarrow \quad 
B = \frac{v_0 + \alpha x_0}{\omega_d}
\]
Now express the solution in phase-amplitude form:
\[
x(t) = C e^{-\alpha t} \cos(\omega_d t - \phi)
\]
The relationship between the constants is:
\[
A = C \cos(\phi), \quad B = C \sin(\phi)
\quad \Rightarrow \quad 
C = \sqrt{A^2 + B^2}, \quad \phi = \tan^{-1} \left( \frac{B}{A} \right)
\]
Substituting the values:
\[
C = \sqrt{x_0^2 + \left( \frac{v_0 + \alpha x_0}{\omega_d} \right)^2}, \quad
\phi = \tan^{-1} \left( \frac{ \frac{v_0 + \alpha x_0}{\omega_d} }{x_0} \right)
= \tan^{-1} \left( \frac{x_0 \omega_d}{-v_0 - \alpha x_0} \right)
\]
These parameters \( A \), \( B \), \( C \), and \( \phi \) uniquely determine the motion for the given initial conditions.

\section{Forced Response}

\subsection{Constant Force}

Consider a system subject to a constant external force \( F(t) = F_0 \). The equation of motion becomes:
\[
m \ddot{x} + c \dot{x} + k x = F_0
\]
The general solution is given by:
\[
x(t) = x_h(t) + x_p(t)
\]
where \( x_h(t) \) is the solution of the homogeneous equation (free response), and \( x_p(t) \) is the particular solution due to the forcing term. For a constant input \( F_0 \), the particular solution will be constant. Assume:
\[
x_p(t) = X_0, \quad \dot{x}_p(t) = \ddot{x}_p(t) = 0
\]
Substituting into the equation:
\[
k X_0 = F_0 \quad \Rightarrow \quad X_0 = \frac{F_0}{k}
\]
So the full solution is:
\[
x(t) = \frac{F_0}{k} + x_h(t)
\]
The constant force simply shifts the equilibrium point of the system.

\subsection{Harmonic Forcing}

\paragraph{Undamped case} Now consider a harmonic excitation:
\[
F(t) = F_0 \cos(\Omega t + \beta)
\]
The equation of motion is:
\[
m \ddot{x} + k x = F_0 \cos(\Omega t + \beta)
\]
Assume a particular solution of the form:
\[
x_p(t) = X_0 \cos(\Omega t + \beta)
\]
Then:
\[
\dot{x}_p(t) = -X_0 \Omega \sin(\Omega t + \beta), \quad
\ddot{x}_p(t) = -X_0 \Omega^2 \cos(\Omega t + \beta)
\]
Substituting into the differential equation:
\[
- m X_0 \Omega^2 \cos(\Omega t + \beta) + k X_0 \cos(\Omega t + \beta) = F_0 \cos(\Omega t + \beta)
\]
\[
\Rightarrow \left( -m \Omega^2 + k \right) X_0 = F_0
\quad \Rightarrow \quad
X_0 = \frac{F_0}{k - m \Omega^2} = \frac{F_0}{k \left( 1 - \frac{\Omega^2}{\omega_0^2} \right)}
\]

\begin{itemize}
  \item If \( \Omega < \omega_0 \): the amplitude is positive.
  \item If \( \Omega > \omega_0 \): the amplitude becomes negative (opposite phase).
  \item If \( \Omega = \omega_0 \): denominator becomes zero, the solution becomes unbounded, this is resonance.
\end{itemize}
When \( \Omega = \omega_0 \), the standard ansatz \( x_p(t) = X_0 \cos(\Omega t + \beta) \) is no longer valid. In this case, a particular solution must include a factor of \( t \) to ensure linear independence. A right proposal is:
\[
x_p(t) = X_0 t \sin(\Omega t + \beta)
\]
Compute the derivatives:

\begin{align*}
\dot{x}_p(t) &= X_0 \sin(\Omega t + \beta) + X_0 t \Omega \cos(\Omega t + \beta) \\
\ddot{x}_p(t) &= X_0 \Omega \cos(\Omega t + \beta) 
- X_0 t \Omega^2 \sin(\Omega t + \beta)
+ X_0 \Omega \cos(\Omega t + \beta) \\
&= 2 X_0 \Omega \cos(\Omega t + \beta) 
- X_0 t \Omega^2 \sin(\Omega t + \beta)
\end{align*}

Now substitute into the undamped equation:
\[
m \ddot{x}_p + k x_p = F_0 \cos(\Omega t + \beta)
\]

\begin{align*}
&m \left[ 2 X_0 \Omega \cos(\Omega t + \beta) - X_0 t \Omega^2 \sin(\Omega t + \beta) \right]
+ k X_0 t \sin(\Omega t + \beta) \\
&= 2 m X_0 \Omega \cos(\Omega t + \beta)
+ X_0 t \left( - m \Omega^2 + k \right) \sin(\Omega t + \beta)
\end{align*}

But if \( \Omega = \omega_0 = \sqrt{k/m} \), then \( k = m \Omega^2 \), so:
\[
- m \Omega^2 + k = 0
\quad \Rightarrow \quad
\text{the term in } t \sin(\cdot) \text{ cancels out}
\]
It left with:
\[
2 m X_0 \Omega \cos(\Omega t + \beta) = F_0 \cos(\Omega t + \beta)
\quad \Rightarrow \quad
X_0 = \frac{F_0}{2 m \Omega}
\]
Thus, the particular solution at resonance becomes:
\[
x_p(t) = \frac{F_0}{2 m \Omega} t \sin(\Omega t + \beta)
\]
This solution grows unbounded in amplitude with time: it represents resonance.

\paragraph{Damped Case} When damping is present, the equation becomes:
\[
m \ddot{x} + c \dot{x} + k x = F_0 \cos(\Omega t + \beta)
\]
Assume a steady-state response of the form:
\[
x_p(t) = \Re\left\{ \tilde{X_0} e^{j(\Omega t + \beta)} \right\}
\quad \text{with} \quad \tilde{X_0} \in \mathbb{C}
\]
Substitute into the equation:
\[
(-m \Omega^2 + j c \Omega + k) \tilde{X_0} = F_0
\quad \Rightarrow \quad
\tilde{X_0} = \frac{F_0}{k - m \Omega^2 + j c \Omega}
\]
Let:
\[
H(\Omega) = \frac{1}{k - m \Omega^2 + j c \Omega}
\quad \Rightarrow \quad
\tilde{X_0} = F_0 H(\Omega)
\]
Then:
\[
x_p(t) = |H(\Omega)| F_0 \cos(\Omega t + \beta + \angle H(\Omega))
\]
This gives the steady-state amplitude and phase shift as a function of the frequency. Let us compute:
\[
|H(\Omega)| = \left| \frac{1}{k - m \Omega^2 + j c \Omega} \right|
= \frac{1}{\sqrt{(k - m \Omega^2)^2 + (c \Omega)^2}}
\]
and:
\[
\angle H(\Omega) = \arctan\left( \frac{-c \Omega}{k - m \Omega^2} \right)
\]

\paragraph{Summary of Damped Harmonic Response} The complete solution is:
\[
x(t) = x_h(t) + x_p(t)
\]
where \( x_h(t) \) is the transient component and \( x_p(t) \) the steady-state one. Graphically, the amplitude and phase shift show distinct behaviors:
\begin{itemize}
  \item For low \( \Omega \): quasi-static regime.
  \item For \( \Omega \approx \omega_0 \): resonance zone.
  \item For high \( \Omega \): inertial regime or seismographic zone.
\end{itemize}

\subsection{Periodic Excitation}

When the external excitation \( f(t) \) is periodic but not purely harmonic, it can be expressed as a Fourier series:
\[
f(t) = a_0 + \sum_{n=1}^{\infty} \left[ a_n \cos(n \Omega t) + b_n \sin(n \Omega t) \right]
\]
or equivalently in complex exponential form:
\[
f(t) = \sum_{n=-\infty}^{+\infty} c_n e^{j n \Omega t}
\]
where the Fourier coefficients are defined as:
\[
c_n = \frac{1}{T} \int_0^T f(t) e^{-j n \Omega t}  dt, \quad T = \frac{2\pi}{\Omega}
\]
The linearity of the system allows us to treat each harmonic component independently and sum the individual responses:
\[
x(t) = \sum_{n=-\infty}^{+\infty} H(j n \Omega) c_n e^{j n \Omega t}
\]
where \( H(j n \Omega) \) is the frequency response function evaluated at the harmonic frequency \( n \Omega \), given by:
\[
H(j n \Omega) = \frac{1}{-m (n \Omega)^2 + j c (n \Omega) + k}
= \frac{1}{k \left( 1 - \frac{(n \Omega)^2}{\omega_0^2} + j \frac{2 h n \Omega}{\omega_0} \right)}
\]
This result holds both for underdamped, overdamped, and critically damped systems, and each harmonic response contributes an oscillation at frequency \( n \Omega \) scaled and phase-shifted according to \( H(j n \Omega) \).

\paragraph{Interpretation:}

\begin{itemize}
    \item The system behaves as a filter, attenuating or amplifying different frequency components depending on proximity to resonance \( \omega_0 \).
    \item Resonance may occur for one or more harmonics if \( n \Omega \approx \omega_0 \).
    \item The total steady-state response \( x_p(t) \) will be a quasi-periodic function, reconstructing the shape of the original excitation filtered by the system's dynamics.
\end{itemize}

\paragraph{Summary:} The complete response of a linear 1-DOF system to a periodic but non-harmonic excitation is obtained by:

\begin{enumerate}
    \item Expanding \( f(t) \) into its Fourier components.
    \item Computing the response of the system to each harmonic using the frequency response function.
    \item Summing all individual harmonic responses via superposition.
\end{enumerate}

\section{Graphical Example}

The following figures illustrate the theoretical response behaviors of a linear 1-DOF mechanical system under various conditions, both in free and forced vibration scenarios:

\begin{figure}[!htb]
    \centering
    \includegraphics[width=0.95\linewidth]{00_Images/1_DOF/01.png}
    \caption{Free vibration responses for different damping ratios \( h \). A lower damping ratio leads to slower energy dissipation and higher oscillation amplitude. The critical case (\( h = 1 \)) marks the transition to non-oscillatory behavior.}
\end{figure}

\begin{figure}[!htb]
    \centering
    \includegraphics[width=0.95\linewidth]{00_Images/1_DOF/02.png}
    \caption{Frequency Response Function (FRF) in magnitude and phase for increasing damping ratios. As damping increases, the resonance peak flattens and the phase transition around \( \omega_0 \) becomes smoother.}
\end{figure}

\begin{figure}[!htb]
    \centering
    \includegraphics[width=0.95\linewidth]{00_Images/1_DOF/03.png}
    \caption{Complete response to harmonic excitation for three different excitation frequencies. Each subplot shows the transient (blue), steady-state (red), and total response (green).}
\end{figure}

\begin{figure}[!htb]
    \centering
    \includegraphics[width=0.95\linewidth]{00_Images/1_DOF/04.png}
    \caption{Response to a multi-harmonic periodic forcing. Top: external torque and its harmonic components. Bottom: corresponding system response, where each harmonic contributes to the total oscillation according to the system's frequency response.}
\end{figure}
